%============================================================================%
% Disciplined heterogeneity: residualise-then-LASSO
%
% Prototype outputs from adhoc/phl/prototype/lasso_heterogeneity.R
% Compile:  latexmk -pdf heterogeneity_methodology.tex
%============================================================================%
\documentclass[10pt, aspectratio=169]{beamer}

\usetheme{default}
\usecolortheme{dove}
\setbeamertemplate{navigation symbols}{}
\setbeamerfont{title}{series=\bfseries}
\setbeamerfont{frametitle}{series=\bfseries}
\setbeamercolor{frametitle}{fg=black}

\usepackage{booktabs}
\usepackage{graphicx}
\usepackage{amsmath}
\usepackage[T1]{fontenc}
\usepackage{lmodern}

\graphicspath{{../figures/phl/}}

\AtBeginSection[]{%
  \begin{frame}[plain]
    \vfill\centering{\usebeamerfont{title}\insertsectionhead}\vfill
  \end{frame}
}

\title{Disciplined Heterogeneity in the Rainfall--Mortality Response}
\subtitle{Allowing many margins, then selecting among them}
\date{\today}

\begin{document}

\begin{frame}[plain]\titlepage\end{frame}

%============================================================================%
\section{The Problem}
%============================================================================%

\begin{frame}{Two mechanisms, one model}
  We want a specification that lets \textbf{both} narratives exist and be
  estimated on the full sample:
  \begin{itemize}
    \item \textbf{Extreme rainfall} damages --- flooding, trauma, waterborne
          disease. Operates at a \emph{daily} scale, days-to-weeks lags,
          concentrated where drainage and density make flooding lethal.
    \item \textbf{Drought} damages --- water scarcity, lost agricultural
          income, malnutrition. Operates on \emph{accumulated deficits} over
          months, with lags of 6--18 months, concentrated where water
          infrastructure is poor and livelihoods are rain-dependent.
  \end{itemize}
  \vspace{0.5em}
  A single smooth $f(\text{daily rainfall})$ \textbf{cannot represent both},
  because they live at different temporal scales. Worse, several channels are
  non-monotone in opposite directions (diarrhoeal disease is U-shaped in
  rainfall; vector-borne disease is hump-shaped), so an all-cause polynomial
  can average genuine effects toward zero.
\end{frame}

\begin{frame}{And a search problem we have already been burned by}
  Testing many candidate margins one at a time does not work. From the PHL
  analysis:
  \begin{itemize}
    \item Five urban-share bins $\times$ three ages gave \textbf{one}
          significant cell out of fifteen --- against $\approx 0.75$ expected
          by chance.
    \item \textbf{Which} cell was significant \emph{moved when the bins were
          redrawn}: quintiles flagged (88.8--100]; a finer grid flagged
          (90--95] while (95--100] went quiet; the \emph{least} urban bin
          cleared for 65+.
    \item Formal tests of the implied threshold --- top-vs-bottom,
          top-vs-pooled-rest, and a scan over cutoffs from 60\% to 95\% ---
          were significant \textbf{nowhere} (smallest $p = 0.11$ across 18
          tests).
  \end{itemize}
  \vspace{0.4em}
  We need a procedure that can consider many margins \emph{without} inheriting
  the multiple-comparison problem of having looked at all of them.
\end{frame}

%============================================================================%
\section{Specification}
%============================================================================%

\begin{frame}{A model that nests both mechanisms}
  \[
    D_{ait} = \exp\Big(
      \underbrace{f_a(\mathbf{E}_{it})}_{\text{daily extremes}} +
      \underbrace{g_a(\mathbf{S}_{it})}_{\text{accumulated deficit}} +
      \sum_j \big[ f^j_a(\mathbf{E}_{it}) V^{\text{flood}}_{ij}
                 + g^j_a(\mathbf{S}_{it}) V^{\text{dry}}_{ij} \big]
      + \mu + \nu \Big)
  \]
  \begin{itemize}
    \item $\mathbf{E}$: daily-scale exposure --- days above local p95/p99,
          excess rain above threshold, Rx1day / Rx5day.
    \item $\mathbf{S}$: accumulated deficit --- SPI/SPEI at 6 and 12 months,
          log-deviation over rolling windows, consecutive dry days.
    \item \textbf{The two exposures carry different interaction sets.} Forcing
          one vulnerability dimension to serve both is precisely what stops
          the current specification from representing both stories.
  \end{itemize}
  \vspace{0.3em}
  {\scriptsize Additive separability in scale is the key restriction. A cross
  term (drought raising flood vulnerability) is allowed but expected to be
  second-order.}
\end{frame}

\begin{frame}{Selection should do less work than it sounds like}
  The strongest discipline is \textbf{over-identification}, not fit. Each
  mechanism makes four independent, falsifiable predictions:
  \begin{center}
    \small
    \begin{tabular}{lll}
      \toprule
       & Flood mechanism & Drought mechanism \\
      \midrule
      Functional shape & wet tail, daily intensity & accumulated deficit \\
      Timescale        & days--weeks               & 6--18 months \\
      Cause of death   & drowning, injury, leptospirosis & intestinal, malnutrition \\
      Vulnerability    & drainage, density, poverty & water infra., aridity, poverty \\
      \bottomrule
    \end{tabular}
  \end{center}
  \vspace{0.5em}
  A model that \emph{fits} best is weak evidence. A mechanism that correctly
  predicts \textbf{shape, lag, cause and heterogeneity simultaneously} is
  strong evidence, because there are far more testable implications than free
  parameters. Let theory restrict the space; use ML to choose within it.
\end{frame}

%============================================================================%
\section{Method}
%============================================================================%

\begin{frame}{Residualise, then LASSO}
  \begin{enumerate}
    \item \textbf{Build wide.} Every exposure term $\times$ age $\times$ every
          candidate vulnerability variable, in \emph{levels and squares} so
          non-monotone heterogeneity is expressible.
    \item \textbf{Partial out the fixed effects} by Frisch--Waugh
          (\texttt{fixest::demean}). The LASSO never sees the $>$18,000 fixed
          effects; they are absorbed exactly.
    \item \textbf{LASSO the interactions, never the core.}
          \texttt{penalty.factor = 0} on $f_a$ and $g_a$: the objects of
          interest are always in and never shrunk. Only the heterogeneity
          space is penalised.
    \item \textbf{Choose $\lambda$ by cluster-blocked CV} --- folds are whole
          provinces.
    \item \textbf{Estimate} the selected model by \texttt{fepois}, and assess
          stability by bootstrap.
  \end{enumerate}
\end{frame}

\begin{frame}{Three details that decide whether this works}
  \begin{itemize}
    \item \textbf{Cluster-blocked folds are not optional.} i.i.d.\ CV lets the
          same municipality inform training and validation; under spatial and
          serial correlation it systematically \emph{underpenalises} and will
          select noise. Folds must be whole provinces (or the theory-driven
          plug-in $\lambda$ of Belloni et al.).
    \item \textbf{Correlated terms are selected arbitrarily.} Post-FE, the
          linear and quadratic precipitation terms correlate at $0.79$. LASSO
          keeps one essentially at random and swaps between folds. Use
          \emph{group} LASSO so a polynomial block enters or exits together,
          or keep functional form out of the selection entirely.
    \item \textbf{Naive post-LASSO inference is invalid.} Reporting selected
          coefficients with ordinary standard errors gives shrunk estimates
          and $p$-values that ignore the search. Valid options: post-double
          selection (Belloni--Chernozhukov--Hansen), cross-fitted double ML,
          or --- as here --- report \emph{selection frequencies} rather than
          $p$-values.
  \end{itemize}
\end{frame}

\begin{frame}{Discovery and confirmation on different data}
  Regularisation controls overfitting \emph{within} a sample. It does not
  protect against the risk that a margin is specific to one country.
  \begin{center}
    \small
    \begin{tabular}{lll}
      \toprule
      Split & Share & Use \\
      \midrule
      Develop & $\sim$40\% & find candidate margins freely; nothing reportable \\
      Tune    & $\sim$20\% & choose $\lambda$, bin edges, lag lengths \\
      Confirm & $\sim$40\% & estimate once; \textbf{these} are the reported numbers \\
      \bottomrule
    \end{tabular}
  \end{center}
  \vspace{0.5em}
  \textbf{Split by geography, not at random.} Random splitting leaks: the same
  municipality appears in both halves and spatial correlation carries
  information across. With a multi-country panel, develop on one country and
  confirm on the others --- out-of-sample in the sense that matters for the
  claim we want to make.
\end{frame}

%============================================================================%
\section{Prototype: Philippines}
%============================================================================%

\begin{frame}{The prototype design}
  Run on the PHL panel to stress-test the machinery before the multi-country
  data arrives (\texttt{05\_lasso\_heterogeneity.R}).
  \begin{itemize}
    \item \textbf{82,593} municipality $\times$ year $\times$ age observations.
    \item \textbf{7 vulnerability variables}: urban share, coastal share,
          log income, long-run precipitation, long-run temperature, log
          population, settlement class --- each in levels and squares.
    \item \textbf{144 design columns}: 18 core exposure terms (unpenalised)
          $+$ \textbf{126 penalised interactions}.
    \item Fixed effects absorbed by Frisch--Waugh: \textbf{13,041} municipality
          $\times$ age $\times$ era and \textbf{5,187} province $\times$ year
          $\times$ age.
    \item $\lambda$ chosen by 10-fold CV with \textbf{provinces as folds}.
  \end{itemize}
\end{frame}

\begin{frame}{Result: the conservative rule selects nothing}
  \begin{columns}[T]
    \begin{column}{0.52\textwidth}
      \begin{center}
        \includegraphics[width=\textwidth]{lasso_cv.png}
      \end{center}
    \end{column}
    \begin{column}{0.46\textwidth}
      \vspace{1em}
      \begin{itemize}
        \item $\lambda_{\text{1se}} = 18.1$ selects \textbf{0 of 126}
              interactions.
        \item $\lambda_{\min} = 0.143$ selects \textbf{49}.
        \item The two differ by \textbf{two orders of magnitude} --- the CV
              curve is very flat, which is what a weak heterogeneity signal
              looks like.
      \end{itemize}
      \vspace{0.5em}
      Under the standard conservative rule, \textbf{no heterogeneity in the
      PHL rainfall response survives province-blocked cross-validation}. That
      is consistent with everything the bin-by-bin analysis found, and it is
      the honest headline.
    \end{column}
  \end{columns}
\end{frame}

\begin{frame}{Which margins survive at the permissive penalty}
  \begin{columns}[T]
    \begin{column}{0.55\textwidth}
      \begin{center}
        \includegraphics[width=\textwidth]{lasso_stability.png}
      \end{center}
    \end{column}
    \begin{column}{0.43\textwidth}
      \vspace{0.8em}
      Selection frequency over \textbf{40 province-block bootstraps} at
      $\lambda_{\min}$. \textbf{22 of 114} terms are selected in $\ge 90\%$ of
      resamples.
      \vspace{0.5em}

      Ordering by dimension: \textbf{population} and \textbf{urban share}
      first (both at 100\%), then coastal, income, long-run precipitation.
      \vspace{0.5em}

      {\scriptsize \textbf{Read with care.} These frequencies are conditional
      on $\lambda_{\min}$, a deliberately permissive penalty chosen because
      the 1-s.e.\ rule kept nothing. High stability here means ``robustly
      selected \emph{given} a generous penalty'', not ``significant''.}
    \end{column}
  \end{columns}
\end{frame}

\begin{frame}{What the prototype actually teaches}
  \begin{itemize}
    \item \textbf{The machinery works} and scales: 126 interactions and 18,000
          fixed effects handled in one pass, minutes of compute.
    \item \textbf{Population size emerges as a margin we had not tested.} It
          was in no bin-by-bin analysis, and it is selected in every bootstrap.
          It is also correlated with urban share --- so part of what we read as
          an ``urban'' effect may be a \emph{population-size} effect, which has
          different mechanisms (health-system capacity, registration quality,
          small-count noise).
    \item \textbf{The PHL null is now much better established.} Not ``we tried
          some bins and found little'', but ``a penalised search over 126
          candidate margins, disciplined by spatially blocked CV, retains
          nothing under the conservative rule''.
    \item That is a far stronger prior to carry into the multi-country panel,
          where the same procedure has real variation to work with.
  \end{itemize}
\end{frame}

%============================================================================%
\section{Caveats and Next Steps}
%============================================================================%

\begin{frame}{What is approximate here}
  \begin{itemize}
    \item \textbf{FWL is exact for linear models; our estimator is Poisson.}
          Selection currently happens in the linear analogue and estimation in
          the Poisson model. Defensible --- selection need not use the
          estimation likelihood --- but not exact. The clean fix is
          demeaning inside the IRLS loop, or a Poisson LASSO with offsets.
    \item \textbf{Bootstrap stability is not inference.} It answers ``how
          reliably is this term selected'', not ``is this coefficient
          non-zero''. Valid inference needs post-double-selection or
          cross-fitting, which is the next build.
    \item \textbf{Correlated polynomial terms should be grouped.} Currently
          selected individually; group LASSO is the correct object.
    \item \textbf{The heterogeneity space is still small.} Seven variables,
          one exposure scale. The design only becomes interesting once daily
          extremes and accumulated deficits are both present.
  \end{itemize}
\end{frame}

\begin{frame}{Next steps}
  \begin{enumerate}
    \item \textbf{Daily and monthly climate} from the cluster --- without both
          exposure scales the two-mechanism model cannot be written down.
    \item \textbf{Cause of death}, which is what separates cancelling
          mechanisms and gives each its own functional form.
    \item \textbf{Water and sanitation coverage}, the vulnerability margin the
          drought literature says governs the deficit channel.
    \item \textbf{Group LASSO plus cross-fitting}, for valid inference on the
          selected margins.
    \item \textbf{Add Brazil and Mexico}, and adopt the develop/confirm split
          across countries --- develop the specification in one, confirm in
          the others.
  \end{enumerate}
  \vspace{0.4em}
  {\scriptsize Step 5 is what turns this from a selection procedure into
  evidence: a margin that is discovered in one country and confirmed in
  another is a different kind of claim from one that survives a penalty in
  the sample where it was found.}
\end{frame}

\end{document}
